%%%%%%%%%%%%%%%%%%%%%%%%%%%%%%%%%%%%%%%%%%%%%%%%%%%%%%%%%%%%%%%%%%%%%%%%%%%%%%%%%%
\begin{frame}[fragile]\frametitle{}
\begin{center}
{\Large Hands-On Example: Predicting House Price}
\end{center}

{\tiny (Ref: Machine Learning is Fun! - Adam Geitgey)}
\end{frame}

%%%%%%%%%%%%%%%%%%%%%%%%%%%%%%%%%%%%%%%%%%%%%%%%%%%%%%%%%%%
\begin{frame}[fragile]\frametitle{Numerical way}
\begin{center}
\includegraphics[width=0.6\linewidth,keepaspectratio]{housing1}
\end{center}
\begin{center}
\includegraphics[width=0.6\linewidth,keepaspectratio]{housing2}
\end{center}
\end{frame}

%%%%%%%%%%%%%%%%%%%%%%%%%%%%%%%%%%%%%%%%%%%%%%%%%%%%%%%%%%%
\begin{frame}[fragile]\frametitle{Machine Learning Type}
\begin{itemize}
\item This is supervised learning. 
\item Knew how much each house sold for, 
\item So, knew the answer to the problem 
\item Need work backwards to figure out the logic.
\end{itemize}
\end{frame}

%
%%%%%%%%%%%%%%%%%%%%%%%%%%%%%%%%%%%%%%%%%%%%%%%%%%%%%%%%%%%%
%\begin{frame}[fragile]\frametitle{Unsupervised Learning}
%What if you didn't know the sale price for each house? Even if all you know is the size, location, etc of each house, it turns out you can still do some really cool stuff. This is called unsupervised learning.
%\begin{center}
%\includegraphics[width=0.6\linewidth,keepaspectratio]{housing3}
%\end{center}
%
%
%\end{frame}
%
%%%%%%%%%%%%%%%%%%%%%%%%%%%%%%%%%%%%%%%%%%%%%%%%%%%%%%%%%%%%
%\begin{frame}[fragile]\frametitle{Unsupervised Learning}
%	\begin{itemize}
%	\item You could have an algorithm that automatically identified different market segments in your data.
%	\item you could do is automatically identify any outlier houses that were way different than everything else. Maybe those outlier houses are giant mansions and you can focus your best sales people on those areas because they have bigger commissions.
%	\end{itemize}
%\end{frame}

%%%%%%%%%%%%%%%%%%%%%%%%%%%%%%%%%%%%%%%%%%%%%%%%%%%%%%%%%%%
\begin{frame}[fragile]\frametitle{Traditional way}
BUT, how to decide which numbers to put? PRAY!!!

\begin{lstlisting}
def estimate_house_sales_price(num_of_bedrooms, sqft, neighborhood):
	price = 0
	price_per_sqft = 200

 	if neighborhood == "hipsterton": 
 		price_per_sqft = 400
	elif neighborhood == "skid row": 
		price_per_sqft = 100
	
	price = price_per_sqft * sqft
	if num_of_bedrooms == 0: 
		price = price - 20000
	else: 
		price = price + (num_of_bedrooms * 1000)
    
	return price
\end{lstlisting}
\end{frame}

%%%%%%%%%%%%%%%%%%%%%%%%%%%%%%%%%%%%%%%%%%%%%%%%%%%%%%%%%%%
\begin{frame}[fragile]\frametitle{Prayer}
	\begin{itemize}
	\item Wouldn't it be better if computer figures out? 
	\item Treat it as black box
	\item Feed Inputs and outputs
	\item That's it!!
	\end{itemize}
\begin{lstlisting}
def estimate_house_sales_price(num_of_bedrooms, sqft, neighborhood):
    price = <computer, plz do some math for me>
    return price
\end{lstlisting}
\end{frame}


%%%%%%%%%%%%%%%%%%%%%%%%%%%%%%%%%%%%%%%%%%%%%%%%%%%%%%%%%%%
\begin{frame}[fragile]\frametitle{Prayer granted!!}

	\begin{itemize}
	\item Notice the magic numbers
	\item .841, 1231.123, 2.324, and 201.234. 
	\item These are weights. 
	\item Better the weights - better the prediction!
	\item Done!!
	\end{itemize}
	
\begin{lstlisting}
def estimate_house_sales_price(num_of_bedrooms, sqft, neighborhood):
	price = 0
	price += num_of_bedrooms * .841
	price += sqft * 1231.123
	price += neighborhood * 2.324
	price += 201.234
	return price
\end{lstlisting}	
\end{frame}

%%%%%%%%%%%%%%%%%%%%%%%%%%%%%%%%%%%%%%%%%%%%%%%%%%%%%%%%%%%
\begin{frame}[fragile]\frametitle{How to figure out? A dumb way}
Step 1: Start with each weight set to 1.0:
\begin{lstlisting}
def estimate_house_sales_price(num_of_bedrooms, sqft, neighborhood):
	price = 0
	price += num_of_bedrooms * 1.0
	price += sqft * 1.0
	price += neighborhood * 1.0
	price += 1.0
	return price
\end{lstlisting}
\end{frame}

%%%%%%%%%%%%%%%%%%%%%%%%%%%%%%%%%%%%%%%%%%%%%%%%%%%%%%%%%%%
\begin{frame}[fragile]\frametitle{How to figure out? A dumb way}
Step 2: Guess/predict for all houses 
\begin{center}
\includegraphics[width=0.8\linewidth,keepaspectratio]{housing4}
\end{center}
Predictions NOT good, right?
\end{frame}

%%%%%%%%%%%%%%%%%%%%%%%%%%%%%%%%%%%%%%%%%%%%%%%%%%%%%%%%%%%
\begin{frame}[fragile]\frametitle{What to do?}
	\begin{itemize}
	\item Actual \$250,000, but guessed \$178,000
	\item Off by \$72,000 for that single house.
	\item Diff can be positive or negative, so square it
	\item Add squared diffs of all houses.
	\item Total: \$86,123,373. 
	\item Thats whole error in the system.
	\item That's how ``wrong'' your function currently is.
	\end{itemize}
\end{frame}

%%%%%%%%%%%%%%%%%%%%%%%%%%%%%%%%%%%%%%%%%%%%%%%%%%%%%%%%%%%
\begin{frame}[fragile]\frametitle{What Next?}
	\begin{itemize}
	\item Average per house error is ``cost''.
	\item Get cost to be zero by playing with the weights. 
	\item Thats the Goal!!!
	\end{itemize}
\end{frame}

%%%%%%%%%%%%%%%%%%%%%%%%%%%%%%%%%%%%%%%%%%%%%%%%%%%%%%%%%%%
\begin{frame}[fragile]\frametitle{Mathematically}
\begin{center}
\includegraphics[width=0.8\linewidth,keepaspectratio]{costeqn}
\end{center}
\end{frame}

%%%%%%%%%%%%%%%%%%%%%%%%%%%%%%%%%%%%%%%%%%%%%%%%%%%%%%%%%%%
\begin{frame}[fragile]\frametitle{Graphically}
Plotting cost values,  for all possible ranges of weights for number\_of\_bedrooms and sqft:
\begin{center}
\includegraphics[width=0.8\linewidth,keepaspectratio]{costgrph}
\end{center}
\end{frame}

%%%%%%%%%%%%%%%%%%%%%%%%%%%%%%%%%%%%%%%%%%%%%%%%%%%%%%%%%%%
\begin{frame}[fragile]\frametitle{Graphically}
Cost is lowest at lowest point of the surface. Ideal.

Weights at that point are the answers!
\begin{center}
\includegraphics[width=0.8\linewidth,keepaspectratio]{costgrph1}
\end{center}
\end{frame}


%%%%%%%%%%%%%%%%%%%%%%%%%%%%%%%%%%%%%%%%%%%%%%%%%%%%%%%%%%%
\begin{frame}[fragile]\frametitle{How to find the lowest cost point?}
	\begin{itemize}
	\item Start somewhere.
	\item Find direction (slope? Derivative? Partial?)
	\item Derivative: tells us which way is downhill for any given point on our graph. 
	\item Move in slope direction.
	\item Adjust our weights to get to next point
	\item ``walking down hill'' towards the lowest point.
	\item That's gradient descent. 
	\item Scikit Learn does this for you, hushsh!!
	\end{itemize}
\end{frame}

%%%%%%%%%%%%%%%%%%%%%%%%%%%%%%%%%%%%%%%%%%%%%%%%%%%%%%%%%%%
\begin{frame}[fragile]\frametitle{Calculus, anybody?}

			\begin{center}
			\includegraphics[width=0.5\linewidth,keepaspectratio]{wtupdate_algo}
			\end{center}
\end{frame}


%%%%%%%%%%%%%%%%%%%%%%%%%%%%%%%%%%%%%%%%%%%%%%%%%%%%%%%%%%%
\begin{frame}[fragile]\frametitle{Modern Implementation with scikit-learn}
\begin{lstlisting}[language=Python]
from sklearn.linear_model import LinearRegression
from sklearn.model_selection import train_test_split
from sklearn.metrics import mean_squared_error, r2_score
import numpy as np, pandas as pd

df = pd.read_csv('housing_data.csv')
X = df[['bedrooms', 'sqft', 'neighborhood_encoded']]
y = df['price']

X_train, X_test, y_train, y_test = train_test_split(
    X, y, test_size=0.2, random_state=42)

model = LinearRegression()
model.fit(X_train, y_train)

y_pred = model.predict(X_test)
print(f"R^2 Score: {r2_score(y_test, y_pred):.3f}")
print(f"RMSE: {np.sqrt(mean_squared_error(y_test, y_pred)):,.2f}")

print(f"Coefficients: {model.coef_}")
print(f"Intercept: {model.intercept_}")
\end{lstlisting}
\end{frame}

% %%%%%%%%%%%%%%%%%%%%%%%%%%%%%%%%%%%%%%%%%%%%%%%%%%%%%%%%%%%
% \begin{frame}[fragile]\frametitle{Feature Engineering for Better Predictions}
% \begin{lstlisting}[language=Python]
% # Create new features
% df['price_per_sqft'] = df['price'] / df['sqft']
% df['is_large'] = (df['bedrooms'] > 3).astype(int)
% df['age'] = 2025 - df['year_built']

% # Polynomial features
% from sklearn.preprocessing import PolynomialFeatures
% poly = PolynomialFeatures(degree=2, include_bias=False)
% X_poly = poly.fit_transform(X)

% # One-hot encoding for neighborhoods
% X_encoded = pd.get_dummies(df, columns=['neighborhood'], 
                           % prefix='neighborhood')
% \end{lstlisting}
% \end{frame}

% %%%%%%%%%%%%%%%%%%%%%%%%%%%%%%%%%%%%%%%%%%%%%%%%%%%%%%%%%%%
% \begin{frame}[fragile]\frametitle{Cross-Validation: More Robust Evaluation}
% \begin{lstlisting}[language=Python]
% from sklearn.model_selection import cross_val_score

% # 5-fold cross-validation
% scores = cross_val_score(model, X, y, cv=5, 
                        % scoring='neg_mean_squared_error')

% rmse_scores = np.sqrt(-scores)
% print(f"RMSE: {rmse_scores.mean():.2f} (+/- {rmse_scores.std():.2f})")

% # Visualize cross-validation results
% import matplotlib.pyplot as plt
% plt.boxplot(rmse_scores)
% plt.ylabel('RMSE ($)')
% plt.title('Cross-Validation Results')
% \end{lstlisting}
% \end{frame}
% %
% %%%%%%%%%%%%%%%%%%%%%%%%%%%%%%%%%%%%%%%%%%%%%%%%%%%%%%%%%%%%
% %\begin{frame}[fragile]\frametitle{Steps to Solve (recap)}
% %	\begin{itemize}
% %	\item  Feature Selection
% %	\item  Feature Scaling
% %	\item  Model Selection
% %	\item  Parameter Selection
% %	\item  Cost Function
% %	\item  Gradient Descent
% %	\item  Evaluation
% %	\end{itemize}
% %\end{frame}
